\section{An exact Lubin--Tate coefficient window}
\label{sec:log-window}
This section advances the coefficient-depth calculation in the incoming
\emph{Cycle signings and Lubin--Tate} report.  It is a local arithmetic
method transfer, independent of RH or GRH.  The relevant source mechanism is
that a sharp integral jet bound can be attained or obstructed by isolating its
first coefficient; compare the source's
\texttt{New/Derivatives/Sharpness/OneCurve.lean}.  The new statement below is
proved directly.  We do not claim that it resolves the bounded-functions
conjecture, which separate November 2026 preprints already announce; see \cite{Colovic}.

Let $p$ be prime, let $f\ge1$ be an integer, put $q=p^f$ and $a=q-1$,
and let $L(X)\in\mathbb Q_p[[X]]$ be the
normalized Lubin--Tate logarithm with
\[
 L(X)=X+O(X^2),\qquad L(pX+X^q)=pL(X).
\]
Normalize $v_p(p)=1$.  Put $v_p(0)=+\infty$.

\begin{theorem}[Exact logarithmic coefficient window]\label{th:log-window}
For every positive integer $n$, put $N=v_p(n)$.  For every
$1\le j\le p^N$ one has
\[
 \boxed{\quad
 v_p\!\left([X^{n+j(q-1)}]L(X)^n\right)
       =N-v_p(j)-j.\quad}
\]
The coefficients outside degrees congruent to $n$ modulo $q-1$ vanish.
Consequently the first nonintegral coefficient after the leading coefficient
occurs exactly in degree
\[
 n+(q-1)J_p(N),\qquad
 J_p(N)=\min\{j\ge1:j+v_p(j)>N\}.
\]
\end{theorem}
\begin{proof}
Choose $\eta$ with $\eta^{q-1}=p$ and write
\[
 W(X)=\eta^{-1}L(\eta X)
     =X\sum_{j\ge0}\lambda_j X^{j(q-1)}.
\]
The functional equation becomes $W(pX+pX^q)=pW(X)$.
Coefficient comparison gives $\lambda_0=1$ and, for $j\ge1$,
\[
 (1-p^{j(q-1)})\lambda_j
   =\sum_{i=0}^{j-1}\lambda_i p^{i(q-1)}
         \binom{1+i(q-1)}{j-i}.
\]
Every denominator on the left is a unit.  Induction proves
$\lambda_j\in\mathbb Z_p$,
\[
 \lambda_1=(1-p^{q-1})^{-1},\qquad
 \lambda_j\in p^{q-1}\mathbb Z_p\quad(j\ge2).
\]
The support assertion follows either from the same recurrence or from the
$\mu_{q-1}$ symmetry of the functional equation.

Let $C(Y)=\sum_j\lambda_jY^j$ and $D(Y)=C(Y)-1$.  There is
$B(Y)\in Y\mathbb Z_p[[Y]]$ such that
\[
 D(Y)=\lambda_1Y(1+p^{q-1}B(Y)).
\]
For every integer $i\ge1$, all nonconstant coefficients of
$(1+p^{q-1}B)^i$ are divisible by $p^{q-1+v_p(i)}$.
Indeed
$v_p\binom{i}{k}\ge v_p(i)-v_p(k)$ and
$(q-1)k-v_p(k)\ge q-1$ for $k\ge1$.
It follows that for $j>i$,
\[
 v_p([Y^j]D(Y)^i)\ge q-1+v_p(i).
\]
Expand $C^n=\sum_i\binom niD^i$.  For $1\le j\le p^N\le n$,
the term $i=j$ has coefficient $\binom nj\lambda_1^j$.  Its valuation is
$N-v_p(j)$: use
$\binom nj=(n/j)\binom{n-1}{j-1}$ and note that the last binomial is a
$p$-adic unit since the lowest $N$ base-$p$ digits of $n-1$ are all $p-1$.
Every term $i<j$ has valuation at least $N+q-1$, by the preceding inequality
and $v_p\binom ni\ge N-v_p(i)$.  Thus no cancellation is possible and
\[
 v_p([Y^j]C^n)=N-v_p(j).
\]
Since $[X^{n+ja}]W^n=p^j[X^{n+ja}]L^n$, the asserted equality follows.
Finally $J_p(N)\le N+1\le p^N$ (also when $N=0$), so the proven range includes
the first failure of integrality.
\end{proof}

The case $p=q=2$ gives $L(X)=\log(1+X)$; the theorem then recovers an
exact coefficient-valuation formula for powers of the usual formal
logarithm.  The proof above uses only $q-1\ge1$, so the stated uniformity
in the residue degree requires no additional hypothesis.

For computation, the threshold also has the closed finite expression
\[
 J_p(N)=\min_{0\le b\le\lfloor\log_p(N+1)\rfloor}
 p^b\left(\left\lfloor\frac{N-b}{p^b}\right\rfloor+1\right).
\]
To see this, the expression for a fixed $b$ is the first positive multiple
of $p^b$ whose sum with $b$ exceeds $N$.  It always satisfies the required
inequality, and every minimizing $j$ is included by taking $b=v_p(j)$.

\subsection{Consequence for the genuine coefficient-depth invariant}
For this application specialize again to $q=p^2$.  Let $K/\mathbb Q_p$
be unramified quadratic, let
$M=\mathcal O_K[[X]]^{\psi=0}$, and write
\[
 f([T](X))=\sum_{m\ge0}c_{f,m}(T)X^m.
\]
Berger--Sprang's $s_0(n)$ is the least $m\ge n$ for which the
$\mathcal O_K$-span of the $c_{f,r}$ with $r\le m$ contains a polynomial of
\emph{degree $n$} and leading coefficient of valuation $-w_q(n)$, where
$w_q(n)=\sum_{r\ge1}\lfloor n/q^r\rfloor$.

\begin{corollary}[A sharp coefficient obstruction]\label{cor:depth-window}
Let $n>0$. If $s_0(n)>n$, then
\[
 s_0(n)\ \ge\ n+(q-1)J_p(v_p(n)).
\]
\end{corollary}
\begin{proof}
Write $E=L^{-1}$.  The coefficient identity
\[
 [T^n]c_{f,m}(T)=[Z^n]f(E(Z))\,[X^m]L(X)^n
\]
follows by substituting $[T](X)=E(TL(X))$.
At $m=n$, the second factor is one.  Since each $c_{f,n}$ is integer-valued
and has degree at most $n$, its leading coefficient has valuation at least
$-w_q(n)$.  Failure at depth $n$, and the integral value group of $K$, imply
that every first factor has valuation at least $1-w_q(n)$.  The theorem
shows that all second factors before the indicated threshold are integral.
Thus all $T^n$ coefficients in the entire generated module have valuation
at least $1-w_q(n)$, which excludes the required degree-$n$ polynomial.
\end{proof}

The equality in Theorem~\ref{th:log-window} resolves the precise first
integrality threshold which the earlier report bounded using
$j+\lfloor\log_pj\rfloor$.  It does \emph{not} prove that $s_0(n)$ equals the
new lower bound.  At later depths a $T^n$ coefficient of the right valuation
may occur in a polynomial of higher degree.  Those higher terms must be
eliminated by integral combinations of genuine trace-kernel rows.  This
is a distinct arithmetic constraint.

\subsection{A necessary correction to the numerical baseline}
Berger's primary-source errata, dated 13 October 2026, identifies an error in
the paper's Appendix A implementation: at each stage it inserts $T^s$
without establishing membership in the module and omits the genuine
possible nonunit diagonal generator.  Consequently its printed experiment
addresses $\operatorname{Pol}(M)+\mathcal O_K[T]$ rather than the defined
$s_0(n)$.  The theoretical definition and the trace-kernel generators remain
usable.  The incoming report's independent proof of the small-digit-sum
criterion is not invalidated by this code error.

The accompanying \texttt{coefficient\_depth.py} starts only from the exact
trace-kernel generators
\[
 X^i(pX+X^q)^j\quad(1\le i\le q-2),\qquad
 (pX^{q-1}+q-1)(pX+X^q)^j.
\]
It forms the logarithm's Carleman matrix $D$, inverts it over the rational
numbers, and forms genuine rows using $D^{-1}\operatorname{diag}(1,T,\dots)D$.
For a cutoff $N$, it suffices to use $j\le N$ in the second family and
$i+j\le N$ in the first: the omitted generators have $X$-adic order
greater than $N$, since $pX+X^q$ has $X$-adic order one.  Every trace-kernel
series is a convergent $\mathcal O_K$-linear combination of these generator
series, and its coefficients through $X^N$ depend only on the finite list.
All rows $c_{f,m}$ with $m\le N$ are retained, including diagonal rows with
nonunit coefficients.  This proves completeness at each finite cutoff.

A descending-degree elimination over $\mathbb Z_{(p)}$ replaces a pivot only
when a smaller valuation is found.  Every row operation is invertible over
that ring.  An echelon row of leading degree $n$ therefore gives exactly the
leading-coefficient ideal of the intersection with polynomials of degree
at most $n$.  Extension to $\mathcal O_K$ preserves this calculation.
No finite $p$-adic precision or added monomial is used.

An independent integer Hermite-normal-form calculation confirms the following
exact values.  The two valuation columns give the leading-coefficient ideal
for degree $n$ after imposing that all higher-degree terms vanish.
\begin{center}
\begin{tabular}{rrrrrr}
$p$ & $q$ & $n$ & $s_0(n)$ & valuation at $s_0(n)-1$ & valuation at $s_0(n)$\\
\hline
2&4&3&6&1&0\\
2&4&5&8&0&-1\\
2&4&7&16&0&-1\\
2&4&12&30&-2&-3\\
3&9&8&16&1&0
\end{tabular}
\end{center}
For $p=2,n=12$, Theorem~\ref{th:log-window} first allows a suitable
$T^{12}$ coefficient in a row at depth $18$, since $v_2(12)=2$ and
$J_2(2)=2$.  The exact module calculation gives $s_0(12)=30$.
Indeed, the genuine generator $g(X)=2X^3+3$ already gives
$[T^{12}]c_{g,18}(T)=15811/9815288$, of valuation $-3$.
Thus reaching the right coefficient valuation alone is insufficient;
eliminating all higher terms can require additional rows.

For the independent check, clear denominators in each rational row using
only $p$-adic units and a common factor $p^K$.  Choose a triangular full-rank
source submatrix whose determinant has valuation $M$.  Its adjugate proves
that $p^M\mathbb Z_p^d$ is already in the scaled row module.  Adding $p^M I$
and reducing entries modulo $p^M$ therefore preserves that module.  Exact
integer column Hermite normal form now gives its leading-degree ideals;
subtracting $K$ restores the original valuations.  The script
\texttt{verify\_depth\_hnf.py} records both the obstruction immediately
before each claimed depth and attainment at that depth.  Monotonicity of
the generated modules proves that these are exact minima, rather than
upper bounds from finding one witness.

The attached JSON files record exact finite computations; a null entry
means only that the cutoff was exceeded.  The finite-ring script
\texttt{verify\_log\_window.py} independently checks the displayed universal
valuation formula for several primes and ranges.  These are reproducible
arithmetic checks supporting the written proof, not Lean kernel checks.

\paragraph{Sources and novelty boundary.}
The primary inputs are \cite{BS,BergerErrata}, namely Berger--Sprang, \emph{Integer-valued polynomials and
$p$-adic Fourier theory}, arXiv:2502.18053, Definition 5.1 and Lemma 5.2;
Berger, \emph{Errata for my articles}, Section 19,
\url{https://perso.ens-lyon.fr/laurent.berger/articles/errata.pdf}; and the
incoming \emph{Cycle signings and Lubin--Tate} manuscript, Section
\emph{When is the first coefficient depth sufficient?}.  The exact
coefficient formula and corrected computation here are new to the reviewed
incoming corpus.  No claim of exhaustive literature priority is made.
